\documentclass[10pt]{article}
\usepackage[margin=.72in]{geometry}
\usepackage{amsmath,amssymb,booktabs,graphicx,microtype}
\usepackage{caption}
\title{SAC and JEPA on Fractal Terrain: Short Experimental Report}
\author{}
\date{September 25, 2026}
\begin{document}
\maketitle

\section{Design}
Each terrain is a $512\times512$ periodic pink-noise intensity field $I$ with
spectral exponent $\alpha$; the physical potential is $\phi=|I|^2$. Intensity
is percentile-normalized (1st--99th), and the potential gradient is divided by
its 95th-percentile magnitude. The source task uses $\alpha=3$; transfer changes
only this exponent. There are 64 training terrain seeds and four disjoint
held-out seeds (100001--100004).

Because $I$ is standardized before squaring, changing $\alpha$ changes spatial
correlation and basin geometry, not the intended one-point scale of $\phi$.
Figure~\ref{fig:terrain-hist} confirms this: whole-field potential histograms
remain close, whereas same-seed spatial traces become smoother as $\alpha$ rises.

\begin{figure}[ht]\centering
\includegraphics[width=.96\linewidth]{dumps/sac_noim_analysis/terrain_value_histograms.png}
\caption{Terrain-value reference for the three current exponents, aggregated
over four held-out seeds. Similar marginals result from standardizing intensity
before $\phi=I^2$; offset traces expose the differing spatial scale.}
\label{fig:terrain-hist}
\end{figure}

The point mass has $m=1$, $\Delta t=.02$ with two semi-implicit Euler substeps,
longitudinal/lateral limits $2/1$, potential scale $4$, and linear damping $.6$.
For $a=(a_{\parallel},a_{\perp})\in[-1,1]^2$,
\[
 \dot v=2a_{\parallel}\hat v+a_{\perp}\hat v_\perp-4\nabla\phi/g_{.95}-.6v .
\]
At each step the agent receives a world-frame $32\times32$ local potential
patch (one world unit per pixel) and a six-dimensional state: two velocity
components, speed, mass, and previous two actions. For the replacement
three-arm experiment, every reset starts at the deterministic global minimum
of that terrain's sampled potential grid. Thus a method cannot obtain initial
speed merely by descending from a high-potential start. Heading and initial
speed in $[0,1]$ remain RNG-controlled and are aligned by matched environment
and evaluation seeds. Episodes have 2,000 steps. Reward is distance travelled
minus $.02\times.05\|a\|^2$; hence mean speed is the direct evaluation metric.
Held-out values use a deterministic policy on the four held-out terrain seeds.

\paragraph{Temporal and spatial optimization scales.}
The JEPA prediction is optimized over $K=8$ environment steps: exactly
$8\times.02=.16$ simulated seconds. Its input is $z_t$ plus the eight
two-dimensional replay actions, and its target is the EMA latent at $t+8$.
This local optimization window is much shorter than the $32\times32$ encoder
field of view (32 world units across), the 2,000-step / 40-second episode, and
SAC's effective discounted horizon of $(1-.99)^{-1}\approx100$ steps (2 s).
The critic therefore learns longer bootstrapped consequences than JEPA's local
transition task. No fixed spatial length follows from eight steps because speed
is state-dependent; on flat terrain the damping-limited $2/.6\approx3.33$
speed would traverse roughly .53 world units in .16 s.

\section{Fresh minimum-start comparison}
All earlier random-start curves are retired. The matched 20k protocol uses SAC,
SAC+JEPA, and SAC+JEPA+variance regularization; every reset starts at the global
terrain minimum. These are distinct ablations:
\[
\begin{array}{lll}
\text{SAC:} & c_{\rm JEPA}=0, & c_{\rm reg}=0,\\
\text{SAC+JEPA:} & c_{\rm JEPA}=.025, & c_{\rm reg}=0,\\
\text{SAC+JEPA+variance:} & c_{\rm JEPA}=.025, & c_{\rm reg}=.01.
\end{array}
\]
The latter two retain the same eight-step predictor and critic-JEPA coefficient;
only the latent regularizer is changed.

Two additional arms use the same $c_{\rm JEPA}=.025$ but replace the variance
penalty by SIGReg ($c_{\rm reg}=.005$) or spherical MMD
($c_{\rm reg}=.01$). Their minimum-start source runs are in progress and are
not silently folded into the completed three-arm result below.

\begin{figure}[ht]\centering
\includegraphics[width=.98\linewidth]{dumps/minstart_v2_analysis/three_arm_curves.png}
\caption{Fresh global-minimum-start comparison. The source column includes
all five completed source arms (SAC, JEPA, variance, SIGReg, spherical MMD).
The two transfer columns presently contain the completed first three arms;
SIGReg and spherical-MMD transfers are pending.}
\end{figure}

\section{Terrain walk examples}
Terrain is not a static image-classification input: the policy must maintain a
speed-producing path through local basins and slopes.  The examples below use
shared held-out terrain tiles and start locations, making the geometric effect
of policy choices directly visible.

\begin{figure}[ht]\centering
\includegraphics[width=.94\linewidth,height=.68\textheight,keepaspectratio]{dumps/spectral_random_start_evaluation/random_start_trajectory_panels.png}
\caption{Representative deterministic terrain walks.  Adjacent SAC and
SAC+JEPA trajectories use the same held-out tile and starting condition.}
\end{figure}

\section{Effect of JEPA}
The SAC actor minimizes
\[
L_{\rm SAC}=\mathbb E[\eta\log\pi(a|o)-\min(Q_1(o,a),Q_2(o,a))],
\]
and each critic minimizes TD MSE against a Polyak target critic ($\gamma=.99$,
$\tau=.005$). For V3 JEPA, an online encoder produces $z_t$ and an EMA encoder
produces a stop-gradient target $\bar z_{t+8}$. The predictor receives eight
replayed actions:
\[
L_{\rm JEPA}=1-\cos\left(q([z_t,a_t,\ldots,a_{t+7}]),\operatorname{sg}[\bar z_{t+8}]\right).
\]
For the current ablation, the actor objective is deliberately simpler:
\[
L_{\rm actor}=L_{\rm SAC}+.025L_{\rm JEPA}+c_{\rm reg}L_{\rm reg}.
\]
The implementation also applies a predictive auxiliary loss to the independent
Q encoders. Inverse dynamics and shuffled-action margin are set exactly to zero
and are not optimized. The archived clean-run CSV records the total critic
objective but not a separate critic-auxiliary component, so the loss figure
below reports the exact logged critic objective without inventing a split.

\section{Effect of regularization terms}
The completed clean screen compares variance, SIGReg, and spherical MMD, all
for 20k transitions with zero inverse and margin coefficients. At source
$\alpha=3$, final held-out speeds are $3.036\pm.036$, $3.028\pm.042$, and
$3.030\pm.012$ respectively. Frozen transfer to $\alpha=2$ yields
$2.998\pm.024$, $3.035\pm.053$, and $3.012\pm.038$; transfer to $\alpha=3.5$
yields $2.975\pm.066$, $3.018\pm.062$, and $2.987\pm.027$. These are one-run,
four-held-out-seed standard deviations, not multi-seed confidence intervals.

\begin{figure}[ht]\centering
\includegraphics[width=.98\linewidth]{dumps/sac_noim_analysis/clean_noim_behavior.png}
\caption{Completed clean no-inverse/no-margin behavioural comparison. Shaded
bands show held-out-seed standard deviations.}\label{fig:clean-behavior}
\end{figure}

\begin{figure}[ht]\centering
\includegraphics[width=.98\linewidth]{dumps/sac_noim_analysis/clean_noim_loss_terms.png}
\caption{Per-condition, exact logged optimization terms. Left: SAC actor
term (solid) and total actor loss (dashed); middle: logged critic objective;
right: weighted actor JEPA and logged variance terms. Inverse and action-margin
terms are identically zero; no critic-JEPA component was logged separately.}
\label{fig:clean-losses}
\end{figure}

\section{Encoder correlation and frozen-transfer representations}
Centered linear CKA was computed on a common 600-step deterministic held-out
observation sequence. Within clean JEPA, $Q_1$--$Q_2$ CKA is $.90$ and
actor--critic CKA is $.76$; the SAC critic pair is $.87$, and cross-model actor
CKA is $.49$. CKA measures representation similarity, not transfer quality.
Moreover, the vanilla checkpoint is at 50k and clean JEPA at 20k transitions,
so this matrix is not a controlled training-budget comparison.

\begin{figure}[ht]\centering
\includegraphics[width=.86\linewidth]{dumps/minstart_v2_analysis/three_arm_encoder_cka_blocks.png}
\caption{Fresh three-run CKA decomposed into actor/$Q_1$/$Q_2$ blocks.}\label{fig:cka}
\end{figure}

\paragraph{Temporal CKA (pending).} The present CKA is the completed 20k
snapshot. A full five-arm source rerun with checkpoints every 2k transitions
will append actor--$Q_1$, actor--$Q_2$, and $Q_1$--$Q_2$ CKA trajectories,
computed on one fixed minimum-start held-out observation sequence. This avoids
mistaking separately sampled trajectories for representational drift.

\begin{figure}[ht]\centering
\includegraphics[width=.94\linewidth]{dumps/minstart_v2_analysis/sharing_temporal_cka.png}
\caption{Per-dump CKA already available from the encoder-sharing source sweep:
actor--$Q_1$, actor--$Q_2$, and $Q_1$--$Q_2$ on fixed minimum-start held-out
observations. The fully-shared run becomes unstable; critic-trained sharing
remains high-similarity and behaviourally stable.}\label{fig:sharing-temporal-cka}
\end{figure}

The following six figures show frozen-encoder transfer diagnostics: terrain/path,
then PCA and UMAP for actor, $Q_1$, and $Q_2$ embeddings on the same trajectory.
PCA and UMAP are fitted per panel, so topology within a panel is diagnostic but
absolute axes cannot be compared across panels. This is especially important
because the actor encoder receives policy plus JEPA gradients, while each critic
encoder receives TD (and, when enabled, predictive) gradients; they need not
parameterize the same latent coordinates. UMAP further has arbitrary rotation,
translation, and scale and is a nonlinear neighbourhood visualization, not a
shared coordinate system. The CKA matrix therefore compares the original
activations on identical observations; it is the appropriate encoder-relation
measurement, whereas separate UMAPs show each encoder's own trajectory geometry.

\begin{figure}[p]\centering
\includegraphics[width=.86\linewidth,height=.78\textheight,keepaspectratio]{dumps/sac_noim_variance_010_frozen_a2/analysis/actor_q_representation_panel_sac_000020000_seed100001.png}
\caption{Frozen encoder, transfer to $\alpha=2$, variance regularization.}\label{fig:rep-a2-var}
\end{figure}
\begin{figure}[p]\centering
\includegraphics[width=.86\linewidth,height=.78\textheight,keepaspectratio]{dumps/sac_noim_sigreg_005_frozen_a2/analysis/actor_q_representation_panel_sac_000020000_seed100001.png}
\caption{Frozen encoder, transfer to $\alpha=2$, SIGReg.}\label{fig:rep-a2-sig}
\end{figure}
\begin{figure}[p]\centering
\includegraphics[width=.86\linewidth,height=.78\textheight,keepaspectratio]{dumps/sac_noim_spherical_mmd_010_frozen_a2/analysis/actor_q_representation_panel_sac_000020000_seed100001.png}
\caption{Frozen encoder, transfer to $\alpha=2$, spherical MMD.}\label{fig:rep-a2-sph}
\end{figure}
\begin{figure}[p]\centering
\includegraphics[width=.86\linewidth,height=.78\textheight,keepaspectratio]{dumps/sac_noim_variance_010_frozen_a35/analysis/actor_q_representation_panel_sac_000020000_seed100001.png}
\caption{Frozen encoder, transfer to $\alpha=3.5$, variance regularization.}\label{fig:rep-a35-var}
\end{figure}
\begin{figure}[p]\centering
\includegraphics[width=.86\linewidth,height=.78\textheight,keepaspectratio]{dumps/sac_noim_sigreg_005_frozen_a35/analysis/actor_q_representation_panel_sac_000020000_seed100001.png}
\caption{Frozen encoder, transfer to $\alpha=3.5$, SIGReg.}\label{fig:rep-a35-sig}
\end{figure}
\begin{figure}[p]\centering
\includegraphics[width=.86\linewidth,height=.78\textheight,keepaspectratio]{dumps/sac_noim_spherical_mmd_010_frozen_a35/analysis/actor_q_representation_panel_sac_000020000_seed100001.png}
\caption{Frozen encoder, transfer to $\alpha=3.5$, spherical MMD.}\label{fig:rep-a35-sph}
\end{figure}

\section{Encoder-sharing alternatives}
The completed experiments use \emph{separate} visual encoders: the actor,
$Q_1$, and $Q_2$ each have a backbone. The main alternatives are:
\begin{center}
\begin{tabular}{llll}
\toprule
Architecture & actor encoder & critic encoder & encoder gradients\\
\midrule
Separate (current) & own & own $Q_1,Q_2$ & each loss updates its own encoder\\
Fully shared & shared & shared & actor, critic, and JEPA all update it\\
Shared, critic-trained & shared read-only & shared & TD critic loss only\\
\bottomrule
\end{tabular}
\end{center}
For critic-trained sharing, $z=e(o)$ feeds the critics while the actor receives
$z_\pi=\operatorname{sg}[z]$:
\[
 L_Q\rightarrow e,\qquad L_\pi\not\rightarrow e,\qquad
 Q_i=Q_i(z,a),\quad a\sim\pi(\operatorname{sg}[z]).
\]
This gives the backbone the critic's TD regression signal while preventing
policy-gradient interference. It is a separate ablation: the optional shared
mode currently implemented here instead uses an actor-owned encoder and detaches
features before critic heads ($L_Q\not\rightarrow e$); critic JEPA is mutually
exclusive with that mode.

\paragraph{Completed controlled sweep.}
Fully shared and critic-trained sharing completed source and frozen $\alpha=2$
and $\alpha=3.5$ transfers with 2k checkpoints. Critic-trained sharing remains
stable; fully shared collapses at $\alpha=3.5$.

\section{Verdict}
\begin{figure}[ht]\centering
\includegraphics[width=.98\linewidth]{dumps/minstart_v2_analysis/sharing_curves.png}
\caption{Completed encoder-sharing source and frozen-transfer comparison.}
\end{figure}
\begin{figure}[ht]\centering
\includegraphics[width=.82\linewidth]{dumps/minstart_v2_analysis/freeze_ablation.png}
\caption{Transfer freeze ablation: both encoders frozen versus critic frozen with actor encoder adaptation.}
\end{figure}
\begin{figure}[ht]\centering
\includegraphics[width=.92\linewidth]{dumps/minstart_v2_analysis/five_arm_encoder_cka_blocks.png}
\caption{Completed five-arm source CKA matrix, decomposed into $3\times3$ actor/$Q_1$/$Q_2$ blocks.}
\end{figure}
JEPA alone does not show a consistent gain over SAC: source endpoints occupy a
narrow $3.04$--$3.09$ range, and plain JEPA is near SAC after both shifts.
Regularization is conditional rather than universally beneficial. Spherical MMD
is the best regularizer at $\alpha=3.5$ (3.080 versus SAC 3.010), while SIGReg
and spherical MMD are lower than SAC at $\alpha=2$ (2.958 and 2.908 versus
3.054). Variance is approximately neutral.

The clearest architectural result is encoder sharing: fully shared gradients
are unstable (fully shared reaches only .011 at frozen $\alpha=3.5$), whereas
critic-trained sharing is stable and strongest among the sharing modes (3.092
at $\alpha=2$, 3.059 at $\alpha=3.5$). The preferred tested design is therefore
\[
z=e(o),\qquad Q_1(z,a),Q_2(z,a),\qquad \pi(\operatorname{sg}[z]),
\]
where TD critic loss updates $e$ and the actor consumes stop-gradient features.
Finally, allowing the actor encoder to adapt while critic encoders remain frozen
hurts the harder $\alpha=2$ transfer (2.369 versus 3.035 when both are frozen),
so freezing both encoders is the better tested transfer protocol. These are
single-training-seed results; small endpoint differences need multi-seed repeats.

\end{document}
